% Paste-ready corrections for the theoretical part of the current
% Job_Market_paper_Kozyrev-2.pdf (61-page version, 2026-07-11 15:38 CEST).

% --------------------------------------------------------------------------
% Replace the joint moment display after Assumptions 1--2.
% --------------------------------------------------------------------------

We interpret the moment restrictions conditionally on the information set
available when the individual forecasts are formed. In particular,
\[
 E\!\left[
 \begin{pmatrix}y_t\\f_t\end{pmatrix}
 \middle|\mathcal F_{t-1}
 \right]
 =
 \begin{pmatrix}\mu_t\\\mu_t\iota_m\end{pmatrix},
 \qquad
 \operatorname{Var}\!\left[
 \begin{pmatrix}y_t\\f_t\end{pmatrix}
 \middle|\mathcal F_{t-1}
 \right]
 =
 \begin{pmatrix}
 \sigma_\epsilon^2&0\\
 0&\Sigma_{e,m}
 \end{pmatrix}.
\]
These equations specify conditional first and second moments; they do not by
themselves impose a joint distribution.

% --------------------------------------------------------------------------
% Replace the discussion surrounding Equation (4).
% --------------------------------------------------------------------------

Throughout the population analysis, $w$ is a fixed weight vector. Since
$f_t=\mu_t\iota_m+u_t$ and $y_t=\mu_t+\epsilon_t$,
\[
 w'f_t-y_t=(w'\iota_m-1)\mu_t+w'u_t-\epsilon_t.
\]
The conditional moment restrictions therefore imply
\[
 E[(w'f_t-y_t)^2\mid\mathcal F_{t-1}]
 =
 (w'\iota_m-1)^2\mu_t^2+\sigma_\epsilon^2+w'\Sigma_{e,m}w.
\]
Under $w'\iota_m=1$, the bias term vanishes. When weights are estimated, we
evaluate them on a new observation independent of the training sample.
Conditional on the training sample, the estimated weights are fixed; the
unconditional out-of-sample risk subsequently averages over training-sample
uncertainty.

% --------------------------------------------------------------------------
% Replace the inflation-factor paragraph.
% --------------------------------------------------------------------------

To summarize the estimation-error trade-off, we use the working factor
\[
 I(T,k)=\frac{T-1}{T-k},\qquad 1\le k<T.
\]
For one correctly specified restricted regression, this factor follows from
the standard i.i.d. Gaussian random-design calculation with $k-1$ free
coefficients and an independent evaluation observation. Assumptions 1--2
alone do not imply that calculation. Moreover, subset regressions fitted on
the same training sample have correlated coefficient-estimation errors.
Consequently, multiplying the infinite-bagged population risk by $I(T,k)$ is
a stylized finite-sample criterion, not an exact risk formula for the feasible
bagged estimator.

% --------------------------------------------------------------------------
% Add after Definition 2.
% --------------------------------------------------------------------------

The quantity $Q_k^{\mathrm{bag},G}$ averages only over the random subset
draws while treating every subset weight as its population value. It should
therefore be distinguished from the out-of-sample risk of feasible RSM, in
which the subset weights are estimated from a common training sample.

% --------------------------------------------------------------------------
% Replace the discussion after Proposition 2.
% --------------------------------------------------------------------------

Proposition~2 establishes two endpoints: the RSM family contains equal
weighting at $k=1$ and the population-optimal linear combination at $k=m$.
It does not establish that weights or risks move monotonically between these
endpoints as $k$ increases. Intermediate values of $k$ may nevertheless be
attractive because they balance approximation to the oracle weights against
weight-estimation error. When $k$ is selected from the data, its tuning cost
must also be taken into account; dominance over both endpoints is therefore
an empirical possibility rather than a consequence of Proposition~2.

% --------------------------------------------------------------------------
% Replace the interpretation after Proposition 3.
% --------------------------------------------------------------------------

This equivalence concerns a conditional mean. A forecast receives more than
the equal global weight if and only if its mean weight, across the subsets in
which it appears, exceeds $1/k$. It need not receive more than $1/k$ in every
subset, and a large regression weight should not by itself be interpreted as
a measure of standalone forecast accuracy.

% --------------------------------------------------------------------------
% Recommended finite-sample version of Proposition 5.
% --------------------------------------------------------------------------

\begin{proposition}[Second-order approximation to average-subset risk]
Under the common-loading model, the exact average-subset population risk is
\[
 Q_k^{\mathrm{sub}}=q+E_S\!\left[\frac{1}{A_S}\right].
\]
A second-order Taylor approximation around $E_S[A_S]=k\bar\tau$ gives
\[
 Q_k^{\mathrm{sub}}
 \approx
 q+\frac{1}{k\bar\tau}
 +\frac{m-k}{k^2(m-1)}\frac{v_\tau}{\bar\tau^3}.
\]
Equivalently, defining
\[
 \eta_k=
 \frac{\operatorname{Var}_S(A_S)}{E_S[A_S]^2}
 =\frac{m-k}{k(m-1)}\frac{v_\tau}{\bar\tau^2},
\]
the approximation is
\[
 E_S[1/A_S]\approx\frac{1}{k\bar\tau}(1+\eta_k).
\]
It is expected to be accurate when $\eta_k$ is small and the distribution of
$A_S$ is not strongly skewed. It becomes exact at $k=m$.
\end{proposition}

% If the author prefers a formal asymptotic remainder, append instead:
%
% Consider a sequence with $m\to\infty$, $k\to\infty$, and precision
% profiles satisfying
% $0<\underline\tau\le\tau_{i,m}\le\overline\tau<\infty$ uniformly.
% Under this sequence, the omitted remainder is $O(k^{-3})$. The subsequent
% rate analysis uses the stronger interior-panel condition $k/m\to0$.

% --------------------------------------------------------------------------
% Replace the clipping rule after Equation (25).
% --------------------------------------------------------------------------

Let $K_T=\min\{m,T-1\}$. The feasible integer choice based on the continuous
approximation is
\[
 k^*=\min\{K_T,\max\{1,\operatorname{round}(k_{\mathrm{sub}}^*)\}\}.
\]

% --------------------------------------------------------------------------
% Replace Proposition 8 and make the same changes in Appendix G.
% --------------------------------------------------------------------------

\begin{proposition}[Rate window under the stylized finite-sample criterion]
Let $m=m_T$, set $K_T=\min\{m_T,T-1\}$, and define
\[
 L_{k,T}^{\mathrm{bag}}
 =
 \left(A+R_{k,T}^{\mathrm{bag}}\right)\frac{T-1}{T-k},
 \qquad
 A=\sigma_\epsilon^2+Q^{\mathrm{opt}}>0,
\]
for $1\le k\le K_T$. Let
\[
 k_T^*\in\arg\min_{1\le k\le K_T}L_{k,T}^{\mathrm{bag}}.
\]
Suppose the marginal regularity bounds hold uniformly along the sequence,
$R_{k,T}^{\mathrm{bag}}$ is uniformly bounded, and
\[
 T\to\infty,
 \qquad
 \frac{T^{1/2}}{m_T}\to0.
\]
Then
\[
 k_T^*=\Omega(T^{1/3}),
 \qquad
 k_T^*=O(T^{1/2}).
\]
Thus, under the stipulated criterion and regularity condition, the minimizing
subset size lies between the cube-root and square-root rates.
\end{proposition}

The proposition concerns the stylized criterion above, not the exact
finite-sample risk of feasible finite-$G$ RSM. Establishing the same rate for
estimated overlapping subset regressions would additionally require control
of cross-subset coefficient-estimation covariances and the finite-$G$ subset
draw term.

% --------------------------------------------------------------------------
% Add at the end of Section 4.
% --------------------------------------------------------------------------

\paragraph{Interpretation of the rate results.}
The exact results in this section concern finite-dimensional population
risks. The square-root formula is obtained by minimizing an approximate
finite-sample criterion, while the cube-root--square-root window describes a
sequence of such criteria as $T$ and the forecast pool grow. These results
motivate finite-sample choices of $k$ but do not make $\sqrt T$ the exact
finite-sample optimum. We therefore evaluate $k=\operatorname{round}(\sqrt T)$
separately in the Monte Carlo experiments and empirical application.

% Preamble presentation fix after loading hyperref:
% \hypersetup{hidelinks}
